\documentclass{article}
\usepackage{amsmath,amssymb,hyperref}
\title{Directed shared-index assembly from native cone geometry}
\date{15 November 2026}
\begin{document}\maketitle
\paragraph{Status and purpose.} The binary, finite-operation and global fixed-interface diagram constructions have independent written-proof review. It returns to semantic interaction rather than evaluator coherence. The result derives a comparison in its actual universal-property direction; it does not invert it to obtain a desired synchronization-release law. No novelty is claimed for limit comparison machinery.

\paragraph{Independent input.} Let $\mathcal D_1,\mathcal D_2,\mathcal E$ be native dialgebra bicategories with the selected incoming lax bilimits of a small ordinary shape $J$, created by their base/output cone equivalences. Let
\[
 H:\mathcal D_1\times\mathcal D_2\longrightarrow\mathcal E
\]
be the already constructed assembly pseudofunctor on native squares and modifications. This is distinct from an underlying normal lax process parallel: the pure-boundary exactness lemma is what makes its action on native squares a pseudofunctor. Assume $H$ preserves pseudo native arrows (those with invertible process-square component), as follows for the prescribed held sequential or parallel operations. Do not assume it preserves the selected bilimits or that its underlying state operation preserves shared-index limits.

For pseudodiagrams $X:J\to\mathcal D_1$, $Y:J\to\mathcal D_2$, write their synchronized objects as $L_X,L_Y$ and their chosen cones as $p_j,q_j$. Their product is a lax cone on $(X,Y)$ in the product bicategory. Applying the actual pseudofunctor $H$ yields a lax cone with vertex $H(L_X,L_Y)$ on the pointwise \emph{same-index} diagram $H(X,Y)$. In detail its leg is $H(p_j,q_j)$; its transition cell is $H$ applied to the pair of input cone cells, with the inverse of the appropriate pseudofunctor compositor inserted. The cone equations are the input cone equations and $H$'s coherence.

\paragraph{Constructed native comparison.} Essential surjectivity of the target cone equivalence supplies
\[
 d_{X,Y}:H(L_X,L_Y)\longrightarrow L_{H(X,Y)}
\]
and an invertible cone modification identifying its projections with the just-described legs. This is a native arrow, with its full state tag and process square; it is not merely a comparison of carrier objects. Applying full faithfulness of that cone equivalence to modifications supplies its naturality for admitted diagram transformations. Here an admitted diagram transformation is one whose components and comparison cells actually induce a functor on the chosen cone categories; pseudonatural diagram transformations always suffice. No blanket claim is made for arbitrary lax transformations of diagrams.

When the input and target cone projections are pseudo, $d_{X,Y}$ is pseudo as a native arrow. Indeed, the created-cone proof determines its process-square component from its projections: each projected component is a paste of the invertible process squares of $H(p_j,q_j)$, the target projection, and the chosen cone identification. Full faithfulness of the output-cone equivalence reflects invertibility, so the process square is invertible. The underlying state map of $d_{X,Y}$ need not be an equivalence. In particular this does not prove $d_{X,Y}$ invertible in the native bicategory.

\paragraph{Coherence without preservation.} Suppose the assembly operations have specified associativity and unit equivalences, natural on native squares, with their coherence equations. For three diagrams, project the two composites of $d$ to each object of the final assembled diagram. Their projected legs are related by the supplied assembly associator. Full faithfulness of the final cone equivalence gives the native two-cell between the two comparison arrows, after retaining the assembly associator and the induced reindexing of the final diagrams. The pentagon equation follows by projecting its two pastes: they become the assembly pentagon and pseudofunctor equations, and faithfulness reflects that equation. The same argument gives unit comparison and triangle. For a unit object $U$, the necessary map is the universal cone map $U\to L_{\mathrm{const}\,U}$ obtained from the identity cone; it is not presumed an equivalence for an arbitrary lax-limit shape. Thus the resulting limit/assembly comparison is lax, not strong, in general.

Every assertion about a supplied assembly equation is made only after its component squares are admitted as native cells. In particular a process-level directed interchanger does not automatically become a diagram transformation of native state doctrines: its boundary tags, native square and modification naturality must first be checked. Once they are supplied, the corresponding comparison diagram can be lifted through the same fully faithful cone functors. This distinguishes deriving cone coherence from assuming missing interaction admission.

\paragraph{Finite-operation substitution.} More generally let a finite multi-interface operation $p$ be a pseudofunctor from a finite product of native doctrines to one output doctrine. Require the selected cones in every doctrine that occurs. The same construction gives
\[
 d_p:p(L_{X_1},\ldots,L_{X_n})\longrightarrow L_{p(X_1,\ldots,X_n)}.
\]
For a nullary operation its source is the specified constant object and its target is the cone on the constant diagram; it is still not presumed invertible. For a composite operation $p(q_1,\ldots,q_r)$, compare its direct $d$ with the paste first applying $p(d_{q_1},\ldots,d_{q_r})$ and then $d_p$ at the diagrams $q_i(\vec X)$. After projection both are the same pointwise composite cone, with the given operation associators and pseudofunctor comparisons. Full faithfulness of the final cone equivalence supplies an invertible native two-cell identifying them. Repeating this argument for three levels proves substitution associativity: the two native two-cell pastes have equal projections. The corresponding unit equation is proved in the same way. Operations may depend on repeated input positions; diagonal selection is an actual product pseudofunctor, and the same selected cone is used at each repeated position. This does not assert that arbitrary noncartesian state tensor copies physical state.

Let $\alpha:p\Rightarrow q$ be an \emph{admitted pseudonatural transformation of native operation pseudofunctors}. Its components may be noninvertible native arrows; its pseudonaturality two-cells have the specified invertible strength. Applying its diagram components gives $L\alpha:L_{p\vec X}\to L_{q\vec X}$. There is an invertible native two-cell
\[
 (L\alpha)\,d_p\;\Longrightarrow\;d_q\,\alpha_{(L_{X_i})}.
\]
Indeed, after projection the two sides are identified by the actual pseudonaturality cell of $\alpha$ at the tuple of cone legs. Naturality for modifications between admitted operation transformations follows by applying their modification equations to those legs. Composition and substitution compatibility follow by the same projection argument. An equation between the supplied modifications remains an equation after lifting, because the cone functor is faithful on native two-cells. This does not require $\alpha$'s component arrows to be equivalences, and those arrows are not inverted.

Consequently a supplied finite operation signature and its admitted directed comparisons act compatibly with shared-index synchronization by these specified lax comparisons. This is a direct local-to-assembled statement, not a claim that an arbitrary name such as ``parallel'' supplies the required native operations. In particular the four-process directed interchanger can enter as $\alpha$ only after its actual state permutation, process square and native pseudonaturality have been constructed. In the prescribed pure-boundary construction the check is finite: use the middle-state permutation as tag, the forward process interchanger with its pure frames as square, and permutation pseudonaturality as the candidate native two-cell. Its admission equation is the original interchanger naturality rectangle at the four native process squares, with the pure-boundary comparisons inserted and cancelled only where invertible. The corresponding modification equations are the original interchanger coherence equations with these same frames. If those local equations are not part of the verified input, this paragraph is conditional on that check and does not manufacture it. Once admitted, the lift above integrates that directed comparison with synchronization without a new global synchronization/interchange axiom.

\paragraph{Application to the original completed relative doctrine.} In the module completion of the original tagged pure/process presentation, the fixed-interface state synchronization theorem constructs the required native lax bilimits from explicit normal-lax cone presentations. Fixed-interface output preservation is derived there from separate locally cocontinuous assembly and restriction-module exactness. The pure-boundary lemma constructs $H$ for the actual native parallel assembly, with output external pair $(A\parallel C,B\parallel D)$. Consequently $d_{X,Y}$ applies to two fixed-interface state diagrams over the same $J$ without requiring the state tensor to preserve that shared index.

This is not the earlier product-index equivalence: it compares two separately synchronized inputs with synchronization of their pointwise joint diagram. Nor does it manufacture the opposite comparison. Its direction is
\[
 \text{assembly of separate synchronizations}
 \longrightarrow\text{synchronization of same-index assembly}.
\]
Selective comparison adjoints or exactness may provide additional reverse structure, but are not premises of this forward construction. In an ambient module doctrine the comparison can be relational or have empty fibers; it is not automatically a deterministic total state transport, an executable coupling, or a physical admission guarantee. The actual native comparison is retained instead of identifying its endpoints.

\paragraph{Global fixed-interface diagram system.} Suppose the native doctrines are the hom bicategories $\mathcal K(A,B)$ of the already constructed process tricategory $\mathcal K$, with its actual sequential and directed parallel operations. Fix a small ordinary shape $J$ and a universe in which the required diagram and cone categories exist. Form hom bicategories
\[
 \mathcal K^J(A,B)=\operatorname{Ps}(J,\mathcal K(A,B)),
\]
whose objects are pseudodiagrams, arrows pseudonatural transformations, and two-cells modifications. External objects $A,B$ remain fixed; the $J$ arrows describe native internal comparisons, not changing external interfaces. The composition operation is the pointwise pair of diagrams followed by the original hom-composition pseudofunctor. The identity is the constant diagram at $1_A$. The original associator and unit pseudonatural adjoint equivalences between the appropriate hom-composition pseudofunctors act pointwise; their coherence modifications and equations likewise act pointwise. Explicitly, naturality along a $J$ arrow is the native naturality of the original associator or unitor; its equation for composable $J$ arrows is that same pseudonaturality equation. Thus no strictification of $\mathcal K$ or deletion of its native witnesses is used. The same construction gives pointwise directed parallel and its already admitted operation comparisons.

Choose the selected lax bilimits and representing cone equivalences coherently by their universal-property construction. On a diagram transformation $a:X\Rightarrow Y$, composing the $X$ cone with $a$ gives a $Y$ cone and hence an arrow $L a:L X\to L Y$. On a modification use fullness of the $Y$ cone equivalence. Comparisons for composition and identity of diagram transformations are the unique invertible native two-cells with the prescribed projected identifications. Their equations follow by faithfulness. Therefore each
\[
 \operatorname{Sync}_{J,A,B}:\mathcal K^J(A,B)\to\mathcal K(A,B)
\]
is a pseudofunctor. Choices are retained as part of this construction; no equality of arbitrarily chosen representatives is asserted.

The sequential operation instance of $d_p$ now gives
\[
 \operatorname{Sync}_J X;\operatorname{Sync}_J Y
 \longrightarrow\operatorname{Sync}_J(X;Y),
\]
and the nullary identity instance gives $1_A\to\operatorname{Sync}_J(\mathrm{const}\,1_A)$. These comparisons are pseudonatural on the hom diagram bicategories. Their associativity and unit modifications are the finite-operation substitution comparisons with the original associator/unitors inserted. Their coherence equations follow by projecting to each $J$ component, where they are exactly the original tricategory equations together with pseudofunctor coherence. Parallel gives its corresponding directed comparison. If the four-process interchanger is admitted at the native operation level as specified above, its compatibility is the lifted $\alpha$ comparison, and its stipulated modification equations lift in the same way.

This is an explicit globally assembled synchronization evaluation with noninvertible sequential/unit/parallel comparisons and specified higher coherence. It is not a strong trihomomorphism, and no normality is claimed: the identity comparison can be noninvertible. We do not use a convention-sensitive label for a lax higher functor without checking that label's complete axiom inventory. The stated data and equations, rather than a name, are the result. Crucially this construction handles arbitrary admitted same-index pseudodiagrams; it does not first replace them by independent products of indices. It still fixes external interfaces and does not turn arbitrary incoherent local process data into an admitted diagram.

\paragraph{Feedback boundary.} A fixed-state closure or joining operation that is already a pseudofunctor on the relevant native doctrines similarly has a canonical comparison from its value on a synchronized cone to synchronization of its image diagram, whenever those two cone constructions exist. To conclude that this comparison is invertible requires preservation of the selected cones. For constant-state closure this includes the constant-state cone condition; for joining it includes the actual shared-index behavior of the state tensor. The present construction does not remove those obstructions by calling both directions feedback compatibility. It provides a directed, witness-preserving interaction operation whose compatibility with each additional admitted operation can be tested by its projected cone cells.

\paragraph{Established background and next audit.} The ordinary one-dimensional precedent is the lax monoidal structure of a limit functor, as implemented in the primary \href{https://leanprover.github.io/hex/docs/Mathlib/CategoryTheory/Monoidal/Limits/Basic.html}{mathlib monoidal-limit development}. The present proof writes out the bicategorical native-square version using the previously established cone equivalences; it does not cite the ordinary implementation as proof of the enriched statement. The exact-review targets are projection strength, admitted diagram transformations, the direction of $d$, and the original-completion application. Arbitrary joint transition laws, unrestricted instantaneous feedback and preservation of every weighted bilimit are not claimed.
\end{document}
